This work investigates whether generating a compressed intermediate representation, followed by deterministic expansion, yields valid BPMN 2.0 models more reliably and more economically in tokens than direct XML generation by language models. A domain-specific language with a formally specified EBNF grammar was designed, together with a deterministic transpiler that expands it into XML validated against the official OMG schema. A corpus of 1,021 processes was built from public sources of textual process descriptions, and a seven-billion-parameter model was specialized through supervised fine-tuning with low-rank adaptation over quantized weights. The evaluation followed a protocol frozen prior to any execution, comprising seven experimental arms, 53 reference processes modeled by experts and three repetitions per item, totaling 1,113 generations. The deterministic component produced valid XML for 100\% of the samples and preserved topological equivalence in 99.6\%, and the proposed representation proved six times more compact than the equivalent logical XML. Contrary to the main hypothesis, describing the language through a prompt reduced reliability: frontier models instructed by the grammar produced valid output in 64.8\% of generations, against 98.1\% for direct XML generation by the same models, with statistical significance and replication across two distinct model families. The specialized small model, in turn, reached 86.8\% validity — twenty-two percentage points above the frontier model instructed by the same grammar — emitting a median of 200 tokens against 631, at a training cost below one dollar. Agreement between reference models produced by different experts for the same process was also measured, yielding a structural F1 of 0.1449, the very range in which all evaluated methods operate; it follows that structural metrics identifying nodes by label saturate close to noise and should not be read as percentage of correctness. The conclusion is that the benefit of compact intermediate representations is not realized through in-context instruction, but through model specialization, which makes fine-tuning a condition for the viability of the approach rather than an accessory stage.

\keywords{BPMN; domain-specific language; language models; supervised fine-tuning; token compression.}
